\documentclass{amsart}
% For dealing with references we use the comment environment
\usepackage{verbatim}
\newenvironment{reference}{\comment}{\endcomment}
%\newenvironment{reference}{}{}
\newenvironment{slogan}{\comment}{\endcomment}
\newenvironment{history}{\comment}{\endcomment}

% For commutative diagrams we use Xy-pic
\usepackage[all]{xy}

% We use 2cell for 2-commutative diagrams.
\xyoption{2cell}
\UseAllTwocells

% We use multicol for the list of chapters between chapters
\usepackage{multicol}

% This is generally recommended for better output
\usepackage{lmodern}
\usepackage[T1]{fontenc}

% For cross-file-references

% Package for hypertext links:
\usepackage{hyperref}

% For any local file, say "hello.tex" you want to link to please

% Theorem environments.
%
\theoremstyle{plain}
\newtheorem{theorem}[subsection]{Theorem}
\newtheorem{proposition}[subsection]{Proposition}
\newtheorem{lemma}[subsection]{Lemma}

\theoremstyle{definition}
\newtheorem{definition}[subsection]{Definition}
\newtheorem{example}[subsection]{Example}
\newtheorem{exercise}[subsection]{Exercise}
\newtheorem{situation}[subsection]{Situation}

\theoremstyle{remark}
\newtheorem{remark}[subsection]{Remark}
\newtheorem{remarks}[subsection]{Remarks}

\numberwithin{equation}{subsection}

% Macros
%
\def\lim{\mathop{\mathrm{lim}}\nolimits}
\def\colim{\mathop{\mathrm{colim}}\nolimits}
\def\Spec{\mathop{\mathrm{Spec}}}
\def\Hom{\mathop{\mathrm{Hom}}\nolimits}
\def\Ext{\mathop{\mathrm{Ext}}\nolimits}
\def\SheafHom{\mathop{\mathcal{H}\!\mathit{om}}\nolimits}
\def\SheafExt{\mathop{\mathcal{E}\!\mathit{xt}}\nolimits}
\def\Sch{\mathit{Sch}}
\def\Mor{\mathop{\mathrm{Mor}}\nolimits}
\def\Ob{\mathop{\mathrm{Ob}}\nolimits}
\def\Sh{\mathop{\mathit{Sh}}\nolimits}
\def\NL{\mathop{N\!L}\nolimits}
\def\CH{\mathop{\mathrm{CH}}\nolimits}
\def\proetale{{pro\text{-}\acute{e}tale}}
\def\etale{{\acute{e}tale}}
\def\QCoh{\mathit{QCoh}}
\def\Ker{\mathop{\mathrm{Ker}}}
\def\Im{\mathop{\mathrm{Im}}}
\def\Coker{\mathop{\mathrm{Coker}}}
\def\Coim{\mathop{\mathrm{Coim}}}

% Boxtimes
%
\DeclareMathSymbol{\boxtimes}{\mathbin}{AMSa}{"02}

%
% Macros for moduli stacks/spaces
%
\def\QCohstack{\mathcal{QC}\!\mathit{oh}}
\def\Cohstack{\mathcal{C}\!\mathit{oh}}
\def\Spacesstack{\mathcal{S}\!\mathit{paces}}
\def\Quotfunctor{\mathrm{Quot}}
\def\Hilbfunctor{\mathrm{Hilb}}
\def\Curvesstack{\mathcal{C}\!\mathit{urves}}
\def\Polarizedstack{\mathcal{P}\!\mathit{olarized}}
\def\Complexesstack{\mathcal{C}\!\mathit{omplexes}}
% \Pic is the operator that assigns to X its picard group, usage \Pic(X)
% \Picardstack_{X/B} denotes the Picard stack of X over B
% \Picardfunctor_{X/B} denotes the Picard functor of X over B
\def\Pic{\mathop{\mathrm{Pic}}\nolimits}
\def\Picardstack{\mathcal{P}\!\mathit{ic}}
\def\Picardfunctor{\mathrm{Pic}}
\def\Deformationcategory{\mathcal{D}\!\mathit{ef}}

\setlength{\emergencystretch}{3em}
\begin{document}
\title[Stacks Project, Tag 0BHZ]{1627 --- Kollar\textquoteright{}s four-way local Noetherian classification: Artinian, one-dimensional regular, depth at least two, or a nontrivial finite punctual modification without maximal associated prime}
\maketitle
\noindent Copyright (C) 2005--2025 Johan de Jong. Stacks Project authors. Source: \href{https://stacks.math.columbia.edu/tag/0BHZ}{Tag 0BHZ}.

\noindent Editorial difficulty note (not author text). Difficulty H1: The theorem classifies all local Noetherian rings by the existence of a nontrivial finite punctual alteration eliminating the maximal associated prime. The proof must construct this alteration either by a torsion quotient or an integral fractional element via the determinantal trick, isolate the exceptional DVR case and prove depth-two rigidity. This structural finite-modification argument is above routine Serre-condition checking..

\noindent Editorial provenance note (not author text). History: excerpt from revision \texttt{a0444\allowbreak 6e57e\allowbreak c1fbc\allowbreak 252a8\allowbreak 71afc\allowbreak ec775\allowbreak 2fb28\allowbreak 07b14}, algebra.tex, lines 28922--29006. Reference presentation was adapted; original-author context and core support blocks were retained or added. The mathematical wording is unchanged. Licensed under GNU FDL 1.2 or later; no invariant sections or cover texts. See ../licenses/COPYING. Context blocks reproduce author definitions and supporting results; they are not additional counted problems.

\begin{definition}
\label{definition-regular-local}
Let $(R, \mathfrak m)$ be a Noetherian local ring of dimension $d$.
\begin{enumerate}
\item A {\it system of parameters of $R$} is a sequence of elements
$x_1, \ldots, x_d \in \mathfrak m$ which generates an ideal of
definition of $R$,
\item if there exist $x_1, \ldots, x_d \in \mathfrak m$
such that $\mathfrak m = (x_1, \ldots, x_d)$ then we call
$R$ a {\it regular local ring} and $x_1, \ldots, x_d$ a {\it regular
system of parameters}.
\end{enumerate}
\end{definition}

\begin{definition}
\label{definition-CM}
Let $R$ be a Noetherian local ring.
Let $M$ be a finite $R$-module.
We say $M$ is {\it Cohen-Macaulay}
if $\dim(\text{Supp}(M)) = \text{depth}(M)$.
\end{definition}

\begin{definition}
\label{definition-maximal-CM}
Let $R$ be a Noetherian local ring.
A finite module $M$ over $R$ is called a {\it maximal Cohen-Macaulay}
module if $\text{depth}(M) = \dim(R)$.
\end{definition}

\begin{definition}
\label{definition-regular-sequence}
Let $R$ be a ring. Let $M$ be an $R$-module. A sequence of elements
$f_1, \ldots, f_r$ of $R$ is called an {\it $M$-regular sequence}
if the following conditions hold:
\begin{enumerate}
\item $f_i$ is a nonzerodivisor on
$M/(f_1, \ldots, f_{i - 1})M$
for each $i = 1, \ldots, r$, and
\item the module $M/(f_1, \ldots, f_r)M$ is not zero.
\end{enumerate}
If $I$ is an ideal of $R$ and $f_1, \ldots, f_r \in I$
then we call $f_1, \ldots, f_r$ an {\it $M$-regular sequence
in $I$}. If $M = R$, we call $f_1, \ldots, f_r$ simply a
{\it regular sequence} (in $I$).
\end{definition}

\begin{definition}
\label{definition-depth}
Let $R$ be a ring, and $I \subset R$ an ideal. Let $M$ be a finite $R$-module.
The {\it $I$-depth} of $M$, denoted $\text{depth}_I(M)$, is defined as follows:
\begin{enumerate}
\item if $IM \not = M$, then $\text{depth}_I(M)$ is the supremum in
$\{0, 1, 2, \ldots, \infty\}$ of the lengths of $M$-regular sequences in $I$,
\item if $IM = M$ we set $\text{depth}_I(M) = \infty$.
\end{enumerate}
If $(R, \mathfrak m)$ is local we call $\text{depth}_{\mathfrak m}(M)$ simply
the {\it depth} of $M$.
\end{definition}

\begin{lemma}
\label{lemma-charpoly-module}
Let $R$ be a ring.
Let $M$ be a finite $R$-module.
Let $\varphi : M \to M$ be an endomorphism.
Then there exists a monic polynomial $P \in R[T]$ such that
$P(\varphi) = 0$ as an endomorphism of $M$.
\end{lemma}

\begin{proof}
Choose a surjective $R$-module map $R^{\oplus n} \to M$, given by
$(a_1, \ldots, a_n) \mapsto \sum a_ix_i$ for some generators $x_i \in M$.
Choose $(a_{i1}, \ldots, a_{in}) \in R^{\oplus n}$ such that
$\varphi(x_i) = \sum a_{ij} x_j$. In other words the diagram
$$
\xymatrix{
R^{\oplus n} \ar[d]_A \ar[r] & M \ar[d]^\varphi \\
R^{\oplus n} \ar[r] & M
}
$$
is commutative where $A = (a_{ij})$. By
Lemma \href{https://stacks.math.columbia.edu/tag/00DX}{lemma charpoly (Tag 00DX)}
there exists a monic polynomial $P$ such that $P(A) = 0$.
Then it follows that $P(\varphi) = 0$.
\end{proof}


\begin{lemma}
\label{lemma-ideal-nonzerodivisor}
Let $R$ be a Noetherian local ring with
maximal ideal $\mathfrak m$. Let $I \subset \mathfrak m$
be an ideal. Let $M$ be a finite $R$-module.
The following are equivalent:
\begin{enumerate}
\item There exists an $x \in I$ which is not a zerodivisor on $M$.
\item We have $I \not \subset \mathfrak q$ for all
$\mathfrak q \in \text{Ass}(M)$.
\end{enumerate}
\end{lemma}

\begin{proof}
If there exists a nonzerodivisor $x$ in $I$,
then $x$ clearly cannot be in any associated
prime of $M$. Conversely, suppose $I \not \subset \mathfrak q$
for all $\mathfrak q \in \text{Ass}(M)$. In this case we can
choose $x \in I$, $x \not \in \mathfrak q$ for all
$\mathfrak q \in \text{Ass}(M)$ by Lemmas
\href{https://stacks.math.columbia.edu/tag/00LC}{lemma finite ass (Tag 00LC)} and \href{https://stacks.math.columbia.edu/tag/00DS}{lemma silly (Tag 00DS)}.
By Lemma \href{https://stacks.math.columbia.edu/tag/00LD}{lemma ass zero divisors (Tag 00LD)} the element $x$
is not a zerodivisor on $M$.
\end{proof}


\begin{lemma}
\label{lemma-regular-mcm-free}
Let $R$ be a regular local ring.
Any maximal Cohen-Macaulay module over $R$ is free.
\end{lemma}

\begin{proof}
Let $M$ be a maximal Cohen-Macaulay module over $R$.
Let $x \in \mathfrak m$ be part of a regular sequence
generating $\mathfrak m$. Then $x$ is a nonzerodivisor
on $M$ by Proposition \ref{proposition-CM-module}, and
$M/xM$ is a maximal Cohen-Macaulay module over $R/xR$.
By induction on $\dim(R)$ we see that $M/xM$ is free.
We win by Lemma \ref{lemma-free-mod-x}.
\end{proof}


\begin{lemma}
\label{lemma-regular-ring-CM}
Let $R$ be a regular local ring and let
$x_1, \ldots, x_d$ be a minimal set of generators
for the maximal ideal $\mathfrak m$. Then
$x_1, \ldots, x_d$ is a regular sequence, and
each $R/(x_1, \ldots, x_c)$ is a regular local ring
of dimension $d - c$. In particular $R$ is Cohen-Macaulay.
\end{lemma}

\begin{proof}
Note that $R/x_1R$ is a Noetherian local ring of dimension $\geq d - 1$
by Lemma \href{https://stacks.math.columbia.edu/tag/00KW}{lemma one equation (Tag 00KW)} with $x_2, \ldots, x_d$
generating the maximal ideal. Hence it is a regular local ring by definition.
Since $R$ is a domain by Lemma \ref{lemma-regular-domain}
$x_1$ is a nonzerodivisor.
\end{proof}


\begin{lemma}
\label{lemma-free-mod-x}
Let $R$ be a Noetherian local ring.
Let $x \in \mathfrak m$.
Let $M$ be a finite $R$-module such that
$x$ is a nonzerodivisor on $M$ and
$M/xM$ is free over $R/xR$.
Then $M$ is free over $R$.
\end{lemma}

\begin{proof}
Let $m_1, \ldots, m_r$ be elements of $M$ which map to
a $R/xR$-basis of $M/xM$. By Nakayama's Lemma \href{https://stacks.math.columbia.edu/tag/00DV}{lemma NAK (Tag 00DV)}
$m_1, \ldots, m_r$ generate $M$. If $\sum a_i m_i = 0$
is a relation, then $a_i \in xR$ for all $i$. Hence
$a_i = b_i x$ for some $b_i \in R$. Hence
the kernel $K$ of $R^r \to M$ satisfies $xK = K$
and hence is zero by Nakayama's lemma.
\end{proof}


\begin{lemma}
\label{lemma-depth-in-ses}
Let $R$ be a local Noetherian ring. Let $0 \to N' \to N \to N'' \to 0$
be a short exact sequence of nonzero finite $R$-modules.
\begin{enumerate}
\item
$\text{depth}(N) \geq \min\{\text{depth}(N'), \text{depth}(N'')\}$
\item
$\text{depth}(N'') \geq \min\{\text{depth}(N), \text{depth}(N') - 1\}$
\item
$\text{depth}(N') \geq \min\{\text{depth}(N), \text{depth}(N'') + 1\}$
\end{enumerate}
\end{lemma}

\begin{proof}
Use the characterization of depth using the Ext groups
$\Ext^i(\kappa, N)$, see Lemma \href{https://stacks.math.columbia.edu/tag/00LW}{lemma depth ext (Tag 00LW)},
and use the long exact cohomology sequence
$$
\begin{matrix}
0
\to \Hom_R(\kappa, N')
\to \Hom_R(\kappa, N)
\to \Hom_R(\kappa, N'')
\\
\phantom{0\ }
\to \Ext^1_R(\kappa, N')
\to \Ext^1_R(\kappa, N)
\to \Ext^1_R(\kappa, N'')
\to \ldots
\end{matrix}
$$
from Lemma \href{https://stacks.math.columbia.edu/tag/00LU}{lemma long exact seq ext (Tag 00LU)}.
\end{proof}


\begin{lemma}
\label{lemma-bound-depth}
Let $(R, \mathfrak m)$ be a Noetherian local ring.
Let $M$ be a nonzero finite $R$-module.
Then $\dim(\text{Supp}(M)) \geq \text{depth}(M)$.
\end{lemma}

\begin{proof}
The proof is by induction on $\dim(\text{Supp}(M))$.
If $\dim(\text{Supp}(M)) = 0$, then
$\text{Supp}(M) = \{\mathfrak m\}$, whence $\text{Ass}(M) = \{\mathfrak m\}$
(by Lemmas \href{https://stacks.math.columbia.edu/tag/0586}{lemma ass support (Tag 0586)} and \href{https://stacks.math.columbia.edu/tag/0587}{lemma ass zero (Tag 0587)}), and hence
the depth of $M$ is zero for example by
Lemma \ref{lemma-ideal-nonzerodivisor}.
For the induction step we assume $\dim(\text{Supp}(M)) > 0$.
Let $f_1, \ldots, f_d$ be a sequence of elements of $\mathfrak m$
such that $f_i$ is a nonzerodivisor on $M/(f_1, \ldots, f_{i - 1})M$.
According to Lemma \href{https://stacks.math.columbia.edu/tag/0AUI}{lemma depth weak sequence (Tag 0AUI)} it suffices to prove
$\dim(\text{Supp}(M)) \geq d$. We may assume
$d > 0$ otherwise the lemma holds. By
Lemma \href{https://stacks.math.columbia.edu/tag/0B52}{lemma one equation module (Tag 0B52)}
we have $\dim(\text{Supp}(M/f_1M)) = \dim(\text{Supp}(M)) - 1$.
By induction we conclude $\dim(\text{Supp}(M/f_1M)) \geq d - 1$
as desired.
\end{proof}


\begin{proposition}
\label{proposition-CM-module}
Let $R$ be a Noetherian local ring, with maximal ideal $\mathfrak m$.
Let $M$ be a Cohen-Macaulay module over $R$ whose support has dimension $d$.
Suppose that $g_1, \ldots, g_c$ are elements of
$\mathfrak m$ such that $\dim(\text{Supp}(M/(g_1, \ldots, g_c)M))
= d - c$. Then $g_1, \ldots, g_c$ is an $M$-regular sequence,
and can be extended to a maximal $M$-regular sequence.
\end{proposition}

\begin{proof}
Let $Z = \text{Supp}(M) \subset \Spec(R)$.
By Lemma \href{https://stacks.math.columbia.edu/tag/00KW}{lemma one equation (Tag 00KW)} in the chain
$Z \supset Z \cap V(g_1) \supset \ldots \supset Z \cap V(g_1, \ldots, g_c)$
each step decreases the dimension at most by $1$. Hence by assumption
each step decreases the dimension by exactly $1$ each time. Thus we
may successively apply Lemma \href{https://stacks.math.columbia.edu/tag/00N5}{lemma CM one g (Tag 00N5)} to the modules
$M/(g_1, \ldots, g_i)$ and the element $g_{i + 1}$.

\medskip\noindent
To extend $g_1, \ldots, g_c$ by one element if $c < d$ we simply
choose an element $g_{c + 1} \in \mathfrak m$ which is not
in any of the finitely many minimal primes of $Z \cap V(g_1, \ldots, g_c)$,
using Lemma \href{https://stacks.math.columbia.edu/tag/00DS}{lemma silly (Tag 00DS)}.
\end{proof}


\begin{lemma}
\label{lemma-regular-domain}
Any regular local ring is a domain.
\end{lemma}

\begin{proof}
We will use that $\bigcap \mathfrak m^n = 0$
by Lemma \href{https://stacks.math.columbia.edu/tag/00IP}{lemma intersect powers ideal module zero (Tag 00IP)}.
Let $f, g \in R$ such that $fg = 0$.
Suppose that $f \in \mathfrak m^a$ and
$g \in \mathfrak m^b$, with $a, b$ maximal.
Since $fg = 0 \in \mathfrak m^{a + b + 1}$
we see from the result of Lemma \ref{lemma-regular-graded}
that either $f \in \mathfrak m^{a + 1}$ or
$g \in \mathfrak m^{b + 1}$. Contradiction.
\end{proof}


\begin{lemma}
\label{lemma-regular-graded}
Let $(R, \mathfrak m, \kappa)$ be a regular local ring of dimension $d$.
The graded ring $\bigoplus \mathfrak m^n / \mathfrak m^{n + 1}$
is isomorphic to the graded polynomial algebra
$\kappa[X_1, \ldots, X_d]$.
\end{lemma}

\begin{proof}
Let $x_1, \ldots, x_d$ be a minimal set of generators
for the maximal ideal $\mathfrak m$, see
Definition \ref{definition-regular-local}.
There is a surjection $\kappa[X_1, \ldots, X_d]
\to \bigoplus \mathfrak m^n/\mathfrak m^{n + 1}$,
which maps $X_i$ to the class of $x_i$ in $\mathfrak m/\mathfrak m^2$.
Since $d(R) = d$ by
Proposition \href{https://stacks.math.columbia.edu/tag/00KQ}{proposition dimension (Tag 00KQ)}
we know that the numerical
polynomial $n \mapsto \dim_\kappa \mathfrak m^n/\mathfrak m^{n + 1}$
has degree $d - 1$. By Lemma \href{https://stacks.math.columbia.edu/tag/00K3}{lemma quotient smaller d (Tag 00K3)} we
conclude that the surjection $\kappa[X_1, \ldots, X_d]
\to \bigoplus \mathfrak m^n/\mathfrak m^{n + 1}$ is an isomorphism.
\end{proof}


\begin{lemma}[Koll\'ar]
\label{lemma-hart-serre-loc-thm}
\begin{reference}
This is taken from a forthcoming paper by
J\'anos Koll\'ar entitled ``Variants of normality for Noetherian schemes''.
\end{reference}
Let $(R, \mathfrak m)$ be a local Noetherian ring.
Then exactly one of the following holds:
\begin{enumerate}
\item $(R, \mathfrak m)$ is Artinian,
\item $(R, \mathfrak m)$ is regular of dimension $1$,
\item $\text{depth}(R) \geq 2$, or
\item there exists a finite ring map $R \to R'$ which is not
an isomorphism whose kernel and cokernel are annihilated by a power
of $\mathfrak m$ such that $\mathfrak m$ is not an associated
prime of $R'$ and $R' \not = 0$.
\end{enumerate}
\end{lemma}

\begin{proof}
Observe that $(R, \mathfrak m)$ is not Artinian if and only if
$V(\mathfrak m) \subset \Spec(R)$ is nowhere dense. See
Proposition \href{https://stacks.math.columbia.edu/tag/00KJ}{proposition dimension zero ring (Tag 00KJ)}. We assume this from now on.

\medskip\noindent
Let $J \subset R$ be the largest ideal killed by a power of $\mathfrak m$.
If $J \not = 0$ then $R \to R/J$ shows that $(R, \mathfrak m)$
is as in (4).

\medskip\noindent
Otherwise $J = 0$. In particular $\mathfrak m$ is not an associated prime
of $R$ and we see that there is a nonzerodivisor $x \in \mathfrak m$ by
Lemma \ref{lemma-ideal-nonzerodivisor}. If $\mathfrak m$
is not an associated prime of $R/xR$ then $\text{depth}(R) \geq 2$
by the same lemma. Thus we are left with the case when there is a
$y \in R$, $y \not \in xR$ such that $y \mathfrak m \subset xR$.

\medskip\noindent
If $y \mathfrak m \subset x \mathfrak m$ then we can consider the
map $\varphi : \mathfrak m \to \mathfrak m$, $f \mapsto yf/x$
(well defined as $x$ is a nonzerodivisor). By the determinantal trick
of Lemma \ref{lemma-charpoly-module} there exists a monic
polynomial $P$ with coefficients in $R$ such that $P(\varphi) = 0$.
We conclude that $P(y/x) = 0$ in $R_x$.
Let $R' \subset R_x$ be the ring generated by
$R$ and $y/x$. Then $R \subset R'$ and $R'/R$ is a finite $R$-module
annihilated by a power of $\mathfrak m$. Thus $R$ is as in (4).

\medskip\noindent
Otherwise there is a $t \in \mathfrak m$ such that $y t = u x$
for some unit $u$ of $R$. After replacing $t$ by $u^{-1}t$
we get $yt = x$. In particular $y$ is a nonzerodivisor.
For any $t' \in \mathfrak m$ we have $y t' = x s$ for some $s \in R$.
Thus $y (t' - s t ) = x s - x s = 0$.
Since $y$ is not a zero-divisor this implies that $t' = ts$ and so
$\mathfrak m = (t)$. Thus $(R, \mathfrak m)$ is regular of dimension 1.

\medskip\noindent
The argument so far shows that every $R$ falls into one of the 4 cases.
To finish we have to show that no two of the properties can hold
simultaneously for each of the 6 combinations of properties. We'll
just show that (2) and (3) each exclude (4); the other cases are left
to the reader.

\medskip\noindent
Assume $R$ is regular of dimension $1$ and that $R \to R'$ is a finite
ring map whose kernel and cokernel are annihilated by a power of $\mathfrak m$
and $\mathfrak m$ is not an associated prime of $R'$. Then $R'$
is a finite $R$-module of depth at least $1$ and hence free by
Lemma \ref{lemma-regular-mcm-free}. Since $R \to R'$ is an isomorphism
in the generic point, we conclude that $R'$ must be free of rank $1$
as an $R$-module. Then $R' \to \text{End}_R(R') \cong R$ is an inverse
to the map $R \to R'$. Thus (4) cannot be true.

\medskip\noindent
Assume $R$ has depth $\geq 2$ and that $R \to R'$ is a finite
ring map whose kernel and cokernel are annihilated by a power
of $\mathfrak m$ and $\mathfrak m$ is not an associated prime of $R'$.
Then $R \to R'$ is necessarily injective (as the kernel would have
both depth $0$ and depth $\geq 1$) and the cokernel has depth
at least $1$ (by Lemma \ref{lemma-depth-in-ses}
and the fact that $R'$ has depth $\geq 1$ by assumption)
whence cannot be supported on
$\{\mathfrak m\}$ unless it is $0$ as well. Thus (4) cannot be true.
\end{proof}
\end{document}
